\documentclass[11pt]{article}
\usepackage{microtype,amsmath,amssymb,amsthm,mathtools,booktabs,array}
\usepackage[margin=1.05in]{geometry}
\usepackage[colorlinks=true,linkcolor=blue!45!black,citecolor=blue!45!black,urlcolor=blue!45!black]{hyperref}
\usepackage{xcolor}
\hypersetup{pdftitle={A sharp coefficient-ring obstruction for Gaussian marker generation},pdfsubject={Working mathematical manuscript}}
\newtheorem{theorem}{Theorem}[section]
\newtheorem{proposition}[theorem]{Proposition}
\newtheorem{lemma}[theorem]{Lemma}
\newtheorem{corollary}[theorem]{Corollary}
\theoremstyle{definition}\newtheorem{definition}[theorem]{Definition}
\theoremstyle{remark}\newtheorem{remark}[theorem]{Remark}
\newcommand{\E}{\mathbb E}
\newcommand{\C}{\mathbb C}
\newcommand{\tr}{\operatorname{tr}}
\newcommand{\wt}{\operatorname{wt}}
\newcommand{\G}{\mathsf G}
\DeclareMathOperator{\Span}{span}
\title{A sharp coefficient-ring obstruction\\for Gaussian marker generation}
\author{}\date{17 November 2026: internal proof note}
\begin{document}\maketitle

\section{Why a fieldwise theorem needs a coefficient guard}

Consider the one-variable Gaussian family
\[
 P_b(Y)=Y^2+bY,\qquad
 F_b(t)=\E[e^{tP_b}]
   =(1-2t)^{-1/2}\exp\!\left(\frac{b^2t^2}{2(1-2t)}\right).
\]
The exponential is a formal series. Its coefficients are polynomials
in $b$, and its constant term is one. This note concerns generation
of its marked period module. It is not a Gaussian Moments Conjecture
counterexample: $\E[P_b]=1$, so the phase is not pure.

Let $R$ be any commutative $\mathbb Q$-algebra and let $b\in R$.
Write $D_R=R[t]\langle\partial_t\rangle$, put
\[
 V_R=\sum_{m\ge0}D_R\Phi_{Y^m}\subset R[[t]],\qquad
 N_R^{\rm act}=D_RF_b,
\]
and define the coefficient torsion ideal
\[
 \Gamma_b(R)=\{c\in R:b^Lc=0\text{ for some }L\ge1\}.
\]
The operators differentiate time; the coefficient parameter $b$
is not a differentiated variable.
Let $\delta_R$ be the point module at $t=1/2$: it is free over
$R$ on $e_j=\partial_t^j e$, $j\ge0$, with
\[
 (1-2t)e_0=0,\qquad
 (1-2t)e_j=2j e_{j-1},\qquad\partial_t e_j=e_{j+1}.
\]

\begin{theorem}\label{thm:allrings}
For every such $R,b$, the exact actual-image defect is
\begin{equation}\label{eq:defect}
 V_R/N_R^{\rm act}\simeq
     \bigl(R/(bR+\Gamma_b(R))\bigr)\otimes_R\delta_R.
\end{equation}
Consequently the unit generates the full actual module if and
only if $bR+\Gamma_b(R)=R$. This is equivalent to existence of
a finite witness
\begin{equation}\label{eq:stabilization}
 b^L(1-br)=0\quad\text{for some }L\ge0,\ r\in R.
\end{equation}
Equivalently, $R$ splits by an idempotent into a component where
$b$ is invertible and a component where $b$ is nilpotent.
\end{theorem}

Over a field this condition always holds: a nonzero $b$ is a unit,
and $b=0$ is nilpotent. Over $R=\mathbb Q[b]$, however,
$\Gamma_b(R)=0$ and the quotient is the nonzero module
$\delta_{\mathbb Q}$. The marker
\[
 \Phi_Y=\frac{bt}{1-2t}F_b
\]
does not belong to $D_RF_b$, even though it does after every
field specialization. Thus the nonunit-content boundary of the
base-marker theorem is a real obstruction in the actual image,
not only in an enlarged source presentation.

The same criterion fails for $\mathbb Q[[b]]$, although it holds
for every finite quotient $\mathbb Q[b]/(b^L)$. Agreement on all
finite coefficient jets does not produce a finite polynomial-Weyl
operator over the completed coefficient ring.

\section{The source and its even and odd parts}

Put $s=1-2t$ and use the exact relative source
\[
 M_R=R[t,Y]/D_YR[t,Y],\qquad
 D_Y=\partial_Y-sY+tb,\qquad
 \nabla=\partial_t+Y^2+bY.
\]
This is an $R[t]$-module quotient, not an ordinary ideal quotient.
The two derivatives commute, and $[\nabla,t]=1$. Let
$N_R=D_R[1]\subset M_R$.

We recall the short polynomial generation calculation, including
the exceptional face. Set $u=1-t$, so $2u-s=1$, and put
$E_j=[Y^{2j}]$, $R_j=[Y^{2j+1}]$. Fibre integration by parts gives
\begin{align}
 s[Y^m]&=tb[Y^{m-1}]+(m-1)[Y^{m-2}],\label{eq:wickrec}\\
 uE_{j+1}&=t\nabla E_j+(2j+1)E_j.\label{eq:connectionrec}
\end{align}
Induction gives $s^{2j}E_j\in R[t][1]$ and
$u^{2j-1}E_j\in N_R$ for $j\ge1$. In the latter induction,
differentiating a power of $u$ raises its denominator by at most
one, and \eqref{eq:connectionrec} adds one more; explicitly the
next cleared expression is
\[
 tu\nabla A_j+\bigl(t(2j-1)+(2j+1)u\bigr)A_j
       =u^{2j+1}E_{j+1},\quad A_j=u^{2j-1}E_j.
\]
Expanding $(2u-s)^{4j-2}=1$ supplies a polynomial B\'ezout identity
between $u^{2j-1}$ and $s^{2j}$. Hence every $E_j$ lies in $N_R$.

Write $K_j=sE_{j+1}-(2j+1)E_j=tbR_j$. Then
\begin{equation}\label{eq:content}
 bR_j=\nabla K_j-K_{j+1}-b^2tE_{j+1}\in N_R.
\end{equation}
This is the time commutator removing $t$, not division by $t$.
Thus $b(M_R/N_R)=0$ over every ring.

Reduction modulo $b$ now computes the entire ambient quotient.
The source commutes with this reduction because it is a cokernel
of a polynomial map. At $b=0$ the even and odd parts are stable
under the connection. The even part is generated by $[1]$, and
the odd part has generators $[Y^{2j+1}]$ with
\[
 s[Y^{2j+1}]=2j[Y^{2j-1}],\qquad
 \nabla[Y^{2j+1}]=[Y^{2j+3}].
\]
It is precisely the point module, by the inverse maps
$[Y^{2j+1}]\mapsto e_j$ and $\partial_t^j e\mapsto[Y^{2j+1}]$.
Therefore
\begin{equation}\label{eq:ambientdefect}
 M_R/N_R\simeq(R/bR)\otimes_R\delta_R.
\end{equation}
This is a polynomial source identity valid with nilpotents and
zero divisors; it is not inferred from generic rank.

\section{Exactly which source classes disappear in the Gaussian image}

The Gaussian map $J_R:M_R\to R[[t]]$ takes $[Y^m]$ to
$\Phi_{Y^m}$. After inverting $s$, recurrence
\eqref{eq:wickrec} gives a free rank-one module with basis $[1]$.
The map there is multiplication by the invertible formal series
$F_b$. The map $R[t,s^{-1}]\to R[[t]]$ is injective: multiply by
a common power of $s$ and use injectivity of polynomials in formal
series. Thus
\begin{equation}\label{eq:kernel}
 \ker J_R=\Gamma_s(M_R).
\end{equation}

We determine its ground classes without a flatness assumption.
The operator $D_Y$ is injective on $R[t,Y]$: for a nonzero
polynomial of highest $Y$ degree, its image has highest term
$-sY$ times that coefficient, and $s$ is a nonzerodivisor in
$R[t]$.

Suppose $s[H]=0$. There is a unique polynomial $Q$ with
$sH=D_YQ$. Reducing at $t=1/2$ gives
\[
 \partial_Y\overline Q+(b/2)\overline Q=0.
\]
Every polynomial solution is
\[
 \overline Q=c\sum_{k\ge0}\frac{(-bY/2)^k}{k!},
 \qquad c\in\Gamma_b(R),
\]
where the sum is finite on $c$. Conversely each such $c$ gives
the source ground class
\begin{equation}\label{eq:ground}
 m_c=\sum_{k=0}^{L-1}\frac{(-tb)^k}{k!}\,c[Y^{k+1}],
                  \qquad b^Lc=0.
\end{equation}
Indeed for $H_{L,c}=\sum_{k<L}(-tbY)^kc/k!$,
\begin{equation}\label{eq:torsionwitness}
 D_YH_{L,c}=-sYH_{L,c}.
\end{equation}
Changing $L$ does not change the polynomial once $b^Lc=0$.

This identifies $\ker(s:M_R\to M_R)$ with $\Gamma_b(R)$.
Here is the well-definedness and completeness check. Replacing
$H$ by $H+D_YV$ replaces $Q$ by $Q+sV$, so $\overline Q$
is intrinsic. If $\overline Q=0$, divide $Q$ by $s$; cancelling
the nonzerodivisor $s$ gives $H\in D_YR[t,Y]$. If
$\overline Q$ has constant coefficient $c$, the class $-m_c$
has the same reduction of its witness, and subtraction reduces
to the zero case. The minus sign simply reflects the witness
convention in \eqref{eq:torsionwitness}. In particular the map
$c\mapsto m_c$ is an $R$-linear isomorphism onto $\ker s$.

For clarity, we also give the elementary point-module argument
that completes this ground calculation. On any module where
$s$ acts locally nilpotently and $[\partial_t,s]=-2$, put
$D=-\partial_t/2$. Then $[D,s]=1$. The finite operator
\[
 \pi(v)=\sum_{j\ge0}\frac{D^js^jv}{j!}
\]
projects to $\ker s$. The sums are finite on each vector.
Successively applying $s$ and this projection gives the unique
expansion
\[
 v=\sum_{j\ge0}\frac{(-1)^j}{j!}D^j\pi(s^jv).
\]
The formulas follow by $sD^j=D^js-jD^{j-1}$ and telescoping;
the top surviving power proves uniqueness. Consequently
\begin{equation}\label{eq:fulltorsion}
 \Gamma_s(M_R)\simeq\Gamma_b(R)\otimes_R\delta_R,
 \quad c\otimes\partial_t^je\longmapsto\nabla^jm_c.
\end{equation}
This is the elementary point-support calculation underlying the
usual Weyl-module description; no analytic comparison or
general image-kernel computation is imported.

\section{The exact actual-image quotient}

In the ambient quotient \eqref{eq:ambientdefect}, the ground
class $m_c$ maps to $c[Y]$. Its $k\ge1$ terms in
\eqref{eq:ground} are either even, or have a coefficient
divisible by $b$, and therefore lie in $N_R$.
The image of \eqref{eq:fulltorsion} in
$(R/bR)\otimes\delta_R$ is exactly
\[
 ((\Gamma_b(R)+bR)/bR)\otimes_R\delta_R.
\]
Since $V_R=M_R/\ker J_R$ and the image of $N_R$ is
$N_R^{\rm act}$, taking the quotient proves
\eqref{eq:defect}.

The criterion \eqref{eq:stabilization} follows directly:
$1=br+c$ with $b^Lc=0$ is equivalent to
$b^L(1-br)=0$. If $L\ge1$, put $z=br$ and $e=z^L$.
Then $z^L=z^{L+1}$, so $e^2=e$, $b(re)=e$, and
$b^L(1-e)=0$. Thus $b$ is invertible on $eR$ and nilpotent
on $(1-e)R$. For $L=0$ it is already invertible. The converse
follows by combining the inverse and nilpotence witnesses on
these two components. This proves all equivalences without
Noetherian or reducedness hypotheses.

\section{Finite witnesses on a nilpotent component}

The nilpotent branch has an explicit operator, not just an
existence argument. Suppose $b^L=0$ and put
$\nu=\lceil L/2\rceil$. Write $h=s^{-1/2}$ and define the
polynomial unit
\[
 A_L(s)=\sum_{k=0}^{\nu-1}\frac{(b^2(2-s)/8)^k}{k!}.
\]
The exact Gaussian formula gives
\[
 A_LF_b=h\sum_{j=0}^{\nu-1}\frac{b^{2j}}{8^jj!}s^{-j},
\]
because $b^2t^2/(2s)=b^2/(8s)-b^2/4+b^2s/8$.
For $E=s\partial_s=-(s/2)\partial_t$, the finite projector
\begin{equation}\label{eq:projector}
 \Pi_\nu(E)=\prod_{j=1}^{\nu-1}\frac{E+j+1/2}{j}
 \qquad\text{satisfies}\qquad
 \Pi_\nu(E)(A_LF_b)=h.
\end{equation}
Each term $s^{-j-1/2}$ is an Euler eigenvector, so the product
kills every $j\ge1$ and fixes $j=0$. Its parameter-derivative
order is at most $\nu-1$; an empty product is one.

Moreover
$\partial_t^kh=(2k-1)!!s^{-k}h$, with the $k=0$ factor one.
These rational constants are units. Conditional Wick expansion
puts every $\Phi_{Y^m}$ in $R[t,s^{-1}]h$, while
\eqref{eq:projector} puts $h$ in $D_RF_b$. Hence
\[
 V_R=D_RF_b=R[t,s^{-1}]h\quad\text{if }b^L=0.
\]
On a unit component the coefficient-content packet gives
ambient generation using the supplied inverse of $b$. For
\eqref{eq:stabilization}, the explicit idempotent $e$ combines
these two finite certificates. Thus nilpotence thresholds,
derivative orders, inverse guards and torsion witnesses can all
be retained by a selector.

\section{A universal countercontrol and the remaining Gaussian problem}

Take $R_0=\mathbb Q[b]$. Here $\Gamma_b(R_0)=0$, so
\eqref{eq:kernel} and \eqref{eq:fulltorsion} make the Gaussian
source map injective. Therefore $V_{R_0}\simeq M_{R_0}$ and
\[
 V_{R_0}/D_{R_0}F_b\simeq\delta_{\mathbb Q}\ne0.
\]
This proves the nonmembership assertion for $\Phi_Y$; its
ambient class maps to the nonzero generator of the quotient.
There is also a concrete divisibility warning: every coefficient
of $\Phi_Y$ is divisible by $b$, yet $\Phi_Y\notin bV_{R_0}$.
Indeed $[Y]$ survives in $M_{R_0}/bM_{R_0}$ as the nonzero
odd point class. Coefficientwise divisibility in formal series
is weaker than divisibility inside the actual image.
The universal actual module is even flat over $\mathbb Q[b]$:
it is torsion-free inside the domain $\mathbb Q[b][[t]]$, and
a torsion-free module over a PID is a filtered union of finite
free submodules. This classical flatness fact is also recorded
in the Stacks Project~\cite[Lemma~15.22.11]{Stacks}.

Nevertheless the actual embedded image does not commute with
the specialization $b=0$. The source does commute with this
cokernel base change, and at $b=0$ its odd point module is killed
by the Gaussian functional. Thus the specialization map has
the exact sequence
\[
 0\longrightarrow\delta_{\mathbb Q}
 \longrightarrow V_{R_0}\otimes_{R_0}R_0/(b)
 \longrightarrow V_{b=0}\longrightarrow0.
\]
For every field value $b=c$, the actual image is generated by
the specialized unit. For every $L\ge1$ the nilpotent quotient
$R_0/(b^L)$ has the finite projector above. Neither statement
removes the universal defect, nor the defect over $\mathbb Q[[b]]$.

The distinction matters for the general programme: the field
base-marker theorem remains valid, and its coefficient-content
guard cannot be dropped by checking fields or finite Artin
thickenings. The older all-ring theorems have their own explicit
identities and are unaffected. This example is not pure and
does not decide the remaining two-base/one-fibre purity-to-safety
question, general vacuum primitivity, or regular/wild realization.

All new statements here have written algebraic proofs. Finite
polynomial controls test the torsion and Euler certificates;
they do not establish the universal quotient calculation.
Classical Gaussian integration, point-module algebra and the
flatness background are credited as such. No new convergent
rational-slope mechanism, nonconvergent formal theorem,
priority, independent replay, formal verification or external
review is claimed. Earlier operator packets, torsion witnesses,
specialization thresholds and completed papers remain unchanged.

\begin{thebibliography}{9}
\bibitem{Stacks}
The Stacks Project Authors, \emph{The Stacks Project},
Section~15.22, Lemma~15.22.11,
\href{https://stacks.math.columbia.edu/tag/0549}{Tag 0549}.
\end{thebibliography}
\end{document}
